% Part: first-order-logic
% Chapter: completeness
% Section: lindenbaums-lemma

\documentclass[../../../include/open-logic-section]{subfiles}

\begin{document}

\iftag{FOL}
      {\olfileid[pt]{fol}{com}{lin}}
      {\olfileid[pt]{pl}{com}{lin}}

\olsection{Lema de Lindenbaum}

\begin{explain}
Provamos agora um lema que mostra que qualquer conjunto consistente de
!!{sentence}s está contido em algum conjunto de sentenças que não é apenas
consistente, mas também !!{complete}. A prova funciona adicionando uma
!!{sentence} de cada vez, garantindo a cada passo que o conjunto permanece
consistente. Fazemos isso de modo que, para todo $!A$, ou $!A$ ou $\lnot
!A$ seja adicionada em algum estágio. A união de todos os estágios dessa
construção então contém ou $!A$ ou sua negação~$\lnot !A$ e
é, portanto, completa. Ela também é consistente, já que nos asseguramos, a cada
estágio, de não introduzir uma inconsistência.
\end{explain}

\begin{lem}[Lema de Lindenbaum]
\ollabel{lem:lindenbaum} Todo conjunto consistente
~$\Gamma$ em uma linguagem~$\Lang{L}$ pode ser
estendido a um conjunto consistente !!{complete}~$\Gamma^*$.
\end{lem}

\begin{proof}
Seja $\Gamma$ consistente.  Sejam $!A_0$, $!A_1$,
\dots{} uma enumeração de todas as !!{sentence}s de~$\Lang L$.
Defina $\Gamma_0 = \Gamma$, e
\[
\Gamma_{n+1} =
\begin{cases}
\Gamma_n \cup \{ !A_n \} & \textrm{se $\Gamma_n \cup \{!A_n\}$ é
  consistente;} \\
\Gamma_n \cup \{ \lnot !A_n \} & \textrm{caso contrário.}
\end{cases}
\]
Seja $\Gamma^* = \bigcup_{n \geq 0} \Gamma_n$.

Cada $\Gamma_n$ é consistente: $\Gamma_0$ é consistente por definição.
Se $\Gamma_{n+1} = \Gamma_n \cup \{!A_n\}$, isso é porque este último
é consistente.  Se não é, $\Gamma_{n+1} = \Gamma_n \cup \{\lnot
!A_n\}$. Temos que verificar que $\Gamma_n \cup \{\lnot !A_n\}$ é
consistente. Suponha que não seja. Então \emph{ambos} $\Gamma_n \cup
\{!A_n\}$ e $\Gamma_n \cup \{\lnot !A_n\}$ são inconsistentes.  Isso
significa que $\Gamma_n$ seria inconsistente por
\iftag{FOL}{%
  \tagrefs{prfAX/{fol:axd:prv:prop:provability-exhaustive},
    prfSC/{fol:seq:prv:prop:provability-exhaustive},
    prfND/{fol:ntd:prv:prop:provability-exhaustive},
    prfTab/{fol:tab:prv:prop:provability-exhaustive}}}{%
  \tagrefs{prfAX/{pl:axd:prv:prop:provability-exhaustive},
    prfSC/{pl:seq:prv:prop:provability-exhaustive},
    prfND/{pl:ntd:prv:prop:provability-exhaustive},
    prfTab/{pl:tab:prv:prop:provability-exhaustive}}},
contrariando a hipótese de indução.

Para todo~$n$ e todo $i < n$, $\Gamma_i \subseteq \Gamma_n$. Isso
segue por uma simples indução sobre~$n$. Para $n=0$, não há $i < 0$,
de modo que a afirmação vale automaticamente.  Para o passo indutivo, suponha que ela
seja verdadeira para~$n$. Mostramos que, se $i < n+1$, então $\Gamma_i \subseteq
\Gamma_{n+1}$. Temos $\Gamma_{n+1} = \Gamma_n \cup \{!A_n\}$ ou $=
\Gamma_n \cup \{\lnot !A_n\}$ por construção. Logo $\Gamma_n \subseteq
\Gamma_{n+1}$. Se $i < n+1$, então $\Gamma_i \subseteq \Gamma_n$ por
hipótese indutiva (se $i < n$) ou pelo fato trivial de que $\Gamma_n
\subseteq \Gamma_n$ (se $i = n$). Obtemos que $\Gamma_i \subseteq
\Gamma_{n+1}$ por transitividade de~$\subseteq$.

Disso segue que $\Gamma^*$ é consistente. Eis o porquê: Seja
$\Gamma' \subseteq \Gamma^*$ finito. Cada $!B \in \Gamma'$ também está
em~$\Gamma_i$ para algum~$i$. Seja $n$ o maior destes. Como
$\Gamma_i \subseteq \Gamma_n$ se $i \le n$, todo $!B \in \Gamma'$ também está
$\in \Gamma_n$, isto é, $\Gamma' \subseteq \Gamma_n$, e
$\Gamma_n$~é consistente. Assim, todo subconjunto finito $\Gamma' \subseteq
\Gamma^*$ é consistente. Por \iftag{FOL}{%
  \tagrefs{prfAX/{fol:axd:ptn:prop:proves-compact},
    prfSC/{fol:seq:ptn:prop:proves-compact},
    prfND/{fol:ntd:ptn:prop:proves-compact},
    prfTab/{fol:tab:ptn:prop:proves-compact}}}{%
  \tagrefs{prfAX/{pl:axd:ptn:prop:proves-compact},
    prfSC/{pl:seq:ptn:prop:proves-compact},
    prfND/{pl:ntd:ptn:prop:proves-compact},
    prfTab/{pl:tab:ptn:prop:proves-compact}}}, $\Gamma^*$ é
  consistente.

Toda !!{sentence} de $\Frm[L]$ aparece na lista usada para
definir~$\Gamma^*$. Se $!A_n \notin \Gamma^*$, então isso é porque
$\Gamma_n \cup \{!A_n\}$ era inconsistente.  Mas então $\lnot !A_n
\in \Gamma^*$, de modo que $\Gamma^*$ é !!{complete}.
\end{proof}

\end{document}
